\chapter{Supplementary Experimental Details}
\label{chap:appendix}

This appendix provides the methodological detail needed to interpret the
experiments in Chapters 3--5. It describes the outputs passed between
generation stages and compares the PiPER capability designs. It then
identifies the task libraries and evaluated tasks, explains the treatment of
reviewed results, and records the model identities and cost assumptions.

\section{Study and Design Outputs}
\label{sec:appendix-stage-outputs}

Study converts robot information into a structured description that Design
can use to specify capabilities. Design then connects those capabilities to
task requirements and executable tests.
Table~\ref{tab:appendix-stage-outputs} summarises the output fields used by
the skeleton-assisted pipeline in Chapter 3. The field names identify the
saved records; their interpretation is given alongside them.

\begin{table}[htbp]
  \centering
  \begingroup
  \small
  \setstretch{1.0}
  \setlength{\tabcolsep}{4pt}
  \renewcommand{\arraystretch}{1.18}
  \caption{Principal fields in the Study and Design outputs.}
  \label{tab:appendix-stage-outputs}
  \begin{tabularx}{\textwidth}{@{}>{\raggedright\arraybackslash}p{29mm}
    >{\raggedright\arraybackslash}p{46mm}
    >{\raggedright\arraybackslash}X@{}}
    \toprule
    \textbf{Record} & \textbf{Fields} & \textbf{Interpretation} \\
    \midrule
    Study: robot identity &
    \path{robot_id}, \path{estimated_class}, \path{dof} &
    Identifies the robot, its estimated embodiment class and degrees of freedom. \\
    Study: control structure &
    \path{joints}, \path{actuators} &
    Joint names, types, axes and recorded control ranges; actuator names,
    joint associations and control ranges. \\
    Study: interaction points &
    \path{ee_site}, \path{ee_body}, \path{graspable_bodies}, \path{notes} &
    Identifies the end-effector reference, graspable bodies and relevant
    observations about the model. \\
    Design: capability interface &
    \path{capabilities} &
    Each entry gives an identifier and method name, describes its effect,
    and specifies request fields, preconditions, timing, invariants and failure behaviour. \\
    Design: behavioural requirements &
    \path{criteria} within each capability &
    Defines each metric, unit, comparator and threshold, together with its
    temporal rule, aggregation and source references. \\
    Design: task coverage &
    \path{task_support} &
    Maps task identifiers to capability identifiers and records the
    rationale for each proposed association. \\
    Design: executable tests &
    \path{scenes}, \path{cases} in \path{scene_cases.yaml} &
    Defines test scenes and concrete requests. Each case binds its
    measurements to simulator objects and the corresponding criteria. \\
    \bottomrule
  \end{tabularx}
  \endgroup
\end{table}

The primary records are \path{study.json} and
\path{design/capability_design.json}. The accompanying
\path{design/scene_cases.yaml} specifies initial conditions, execution limits
and measurement bindings. The framework also records the robot configuration
and task-library snapshot associated with the design.
These records preserve the distinction between a capability's intended
effect and the concrete test used to assess it.

Task coverage in \path{task_support} is a design claim. Behavioural validation
assesses the generated implementation against the specified cases, while
downstream evaluation assesses complete tasks.
Scene preparation establishes that a test can execute and obtain its
measurements; it does not establish that the driver will pass.

\section{Variation among PiPER Capability Designs}
\label{sec:appendix-piper-designs}

Table~\ref{tab:piper-design-run01} in Chapter 3 gives a complete worked
example, including test requests and measured results.
Table~\ref{tab:appendix-piper-comparison} retains the differences needed to
interpret all five runs. Every design included end-effector positioning,
wrist rotation, gripper aperture and contact establishment.
The remaining capability provided ordered waypoint motion in runs 01--04
and position holding in run 05.

\begin{table}[htbp]
  \centering
  \begingroup
  \small
  \setstretch{1.0}
  \setlength{\tabcolsep}{3pt}
  \renewcommand{\arraystretch}{1.2}
  \caption{Differences among the five generated PiPER capability designs.}
  \label{tab:appendix-piper-comparison}
  \begin{tabularx}{\textwidth}{@{}>{\centering\arraybackslash}p{8mm}
    >{\raggedright\arraybackslash}p{36mm}
    >{\raggedright\arraybackslash}X@{}}
    \toprule
    \textbf{Run} & \textbf{Fifth capability} &
    \textbf{Contact criterion} \\
    \midrule
    01 & Two ordered waypoints &
    Peak normal force over 1--3~s $\geq 0.5$~N. \\
    02 & Three ordered waypoints &
    Minimum normal force over 1--2~s $\geq 1$~N. \\
    03 & Two ordered waypoints &
    Mean normal force over 1--2~s $\geq 0.5$~N, and terminal
    end-effector error $\leq 30$~mm. \\
    04 & Two ordered waypoints &
    Peak normal force over 1--3~s within $[0.5,60]$~N. \\
    05 & Position hold; peak error over 0.2--1~s $\leq 20$~mm &
    Peak normal force over 0.5--1.5~s $\geq 1$~N. \\
    \bottomrule
  \end{tabularx}
  \par\smallskip
  \noindent\raggedright\emph{Notes:} Ordered waypoint error is bounded by
  30~mm in runs 01--04. Terminal joint7 error is bounded by 3~mm in
  runs 01, 02 and 05, and by 5~mm in runs 03 and 04.
  Run 05 additionally requires
  $|\mathrm{joint8}+\mathrm{joint7}|\leq 3$~mm.
  Contact forces refer to each case's specified geometry pair.
  These are generated test definitions, not measured outcomes.\par
  \endgroup
\end{table}

The contact criteria assess different properties of a trajectory.
A peak criterion can be satisfied by a brief force excursion, whereas
the minimum criterion requires the threshold throughout the recorded window.
A mean criterion can tolerate intervals below its threshold.
The contact results are therefore reported separately in Chapter 3.

Measurement bindings also affected interpretation. Run 02 assigned force
measurement to a visual geometry without collision enabled.
Contact through other gripper geometries could not satisfy that specified
pair's criterion. Complete per-run requests and bindings are retained under
\path{expriment/chapter3_piper/data/runs/}.

\section{Task Libraries and Evaluation Scope}
\label{sec:appendix-tasks}

The task libraries supplied the task requirements considered during Design.
Each relevant robot package contained 20 entries, but downstream experiments
executed selected subsets. An entry's inclusion in a library therefore
does not imply that it was executed or successfully completed.
Table~\ref{tab:appendix-task-libraries} identifies the libraries used for the
robots studied in this thesis and the corresponding evaluation subsets.

\begin{table}[htbp]
  \centering
  \begingroup
  \small
  \setstretch{1.0}
  \setlength{\tabcolsep}{4pt}
  \renewcommand{\arraystretch}{1.18}
  \caption{Task-library inputs and downstream evaluation subsets.}
  \label{tab:appendix-task-libraries}
  \begin{tabularx}{\textwidth}{@{}>{\raggedright\arraybackslash}p{28mm}
    >{\raggedright\arraybackslash}p{40mm}
    >{\raggedright\arraybackslash}X@{}}
    \toprule
    \textbf{Robot / chapter} & \textbf{Package and version} &
    \textbf{Tasks evaluated downstream} \\
    \midrule
    PiPER / 3 & \path{piper/1.0.0} &
    Reaching, pushing and pick-and-place; three spatial layouts. \\
    SO-101 / 4 & \path{robotstudio_so101/1.0.4} &
    Push to Goal, Sweep into Goal, Pick and Place, Pick and Place with Wall,
    and Dial Turn. \\
    Unitree Go2 / 4 & \path{unitree-go2-stock-}\newline
    \path{12dof/1.0.2} &
    Single-step Ascent, Single-step Descent, Mixed-obstacle Path,
    Pause-table Transition and Weave Poles. \\
    Franka Panda / 5 & \path{franka_panda/1.0.0} &
    Reaching, pushing, pick-and-place, drawer opening and dial turning;
    three spatial layouts. \\
    LEAP Hand / 5 & \path{leap_hand/1.0.0} &
    All-fingertip reaching, joint pose matching, object holding,
    block manipulation and valve turning. \\
    Skydio X2 / 5 & \path{skydio_x2/1.0.0} &
    Takeoff and hover, stationary holding, climb--transit--descend return,
    ordered waypoints and point-of-interest orbiting. \\
    \bottomrule
  \end{tabularx}
  \endgroup
\end{table}

\subsection{Catalogue Contents}

The following inventory preserves the scope of the Design inputs without
repeating every task's numerical clauses. The complete definitions are
stored in \path{catalog.json} within the packages listed in
Table~\ref{tab:appendix-task-libraries}, under
\path{AA1/auto_adapter/task_libraries/}. Each package's
\path{sources.json} records the source material and adaptations.

\begingroup
\small
\setstretch{1.12}
\paragraph{Serial-arm manipulation: PiPER, SO-101 and Panda.}
The three packages share the same catalogue: target reaching; pushing;
pick-and-place; pick-and-place with a wall; pushing around a wall;
sweeping; drawer opening and closing; frontal and top-down button pressing;
handle pressing and pulling; door opening and closing; faucet opening;
dial turning; lever pulling; side peg insertion; bin picking; and extraction
from a hole. Sharing task names does not make their robot-specific scenes
or generated capability tests identical.

\paragraph{Unitree Go2.}
Entries T01--T20 comprise eight-direction velocity tracking; step ascent;
step descent; payload step ascent; high-rate turning; a mixed-obstacle path;
crossing-ramp, compliant-floor, pallet-and-pipe, diagonal-rail and
negative-obstacle lanes; block, stair, stepping-stone and pole terrain;
a pause-table transition; weave poles; an A-frame; a broad jump; and the
complete Barkour course.

\paragraph{LEAP Hand.}
The first thirteen entries cover pinch translation, dynamic tripod motion,
squeeze, twiddle, rock, opposed-grasp pivot, radial roll, index roll,
full roll, rotary step, interdigital step, linear step and palmar slide.
The remaining seven cover all-fingertip reaching, block-pose matching,
joint-pose matching, fixed-axis turning, continuous rotation tracking,
object holding and two-ball Baoding circulation.

\paragraph{Skydio X2.}
Entries T01--T20 cover takeoff, landing, stationary holding,
climb--transit--descend return, ordered waypoints, a continuous flight path,
point-of-interest orbiting, moving-bearing tracking, subject-relative escort,
panorama motion, structure approach, obstacle transit, flight beneath a height
bound, survey lanes, tower inspection, concave-structure observation,
under-bridge traversal, lateral baseline acquisition, scan-route resumption
and geofenced entry and exit.
These entries describe aircraft motion and geometric conditions in MuJoCo.
They do not assess onboard perception, image quality or a physical flight SDK.
\endgroup


\subsection{Sources and Adaptations}

Table~\ref{tab:appendix-standard-sources} distinguishes task-source material
from the local implementation of an evaluation.
A source may supply an operation, a metric or a numerical threshold;
these forms of support are not interchangeable.
In particular, a reported policy performance value used to construct a local
threshold is not an independently established acceptance standard.

\begin{table}[htbp]
  \centering
  \begingroup
  \small
  \setstretch{1.0}
  \setlength{\tabcolsep}{4pt}
  \renewcommand{\arraystretch}{1.18}
  \caption{Source material for the retained task libraries.}
  \label{tab:appendix-standard-sources}
  \begin{tabularx}{\textwidth}{@{}>{\raggedright\arraybackslash}p{36mm}
    >{\raggedright\arraybackslash}X@{}}
    \toprule
    \textbf{Task family} & \textbf{Source and use} \\
    \midrule
    Serial-arm manipulation &
    Meta-World supplies task definitions and metric thresholds
    \citep{yu2020metaworld,farama2026metaworld}. The local packages adapt
    robot models, scenes and requests. \\
    LEAP dexterity tasks &
    The Elliott and Connolly benchmark motivates the first thirteen
    operations, adapted into fixed three-trial protocols
    \citep{coulson2021elliott}. \\
    LEAP reaching and object tasks &
    Gymnasium-Robotics, ROBEL and MyoSuite supply the remaining task families
    \citep{delazcano2026gymnasiumrobotics,ahn2020robel,robel2020software,
    caggiano2022myosuite,myohub2026myosuite}.
    Adaptations include morphology-scaled fingertip tolerance and
    robot-specific joint and object bindings. \\
    Go2 motion and terrain &
    Locomotion studies inform the motion, terrain and morphology-scaled
    requirements \citep{lee2020quadrupedal,miki2022perceptive,
    shi2023terrain,hutter2016anymal}. Their reported commands or performance
    values serve as inputs to local task definitions. \\
    Go2 obstacle courses &
    The ICRA Quadruped Robot Challenge and Barkour provide obstacle-course
    requirements \citep{jacoff2023quadruped,caluwaerts2023barkour,
    yu2024barkourdesign}. The catalogue records the adapted geometry,
    completion conditions and timing. \\
    Skydio flight tasks &
    The local catalogue proposes numerical tolerances and scene arrangements.
    These values are uncalibrated local requirements; they are not
    manufacturer accuracy specifications or demonstrated hardware performance. \\
    \bottomrule
  \end{tabularx}
  \endgroup
\end{table}

Robot-description files provide morphology and actuator information.
Task thresholds require a separate justification, and capability validation
remains distinct from downstream task evaluation.

\subsection{Reading the Evaluation Criteria}
\label{sec:appendix-evaluation-criteria}

A task passes only when all of its required clauses are satisfied.
The measurement object and time rule are part of each clause:
an end-effector error differs from an object error, and an early target hit
differs from success at a prescribed endpoint.
Chapter 3 specifies its PiPER task geometry and criteria directly.
For the Chapter 5 cases,
Table~\ref{tab:appendix-task-criteria} records the distinctions needed to
interpret the reported outcomes.

\begin{table}[H]
  \centering
  \begingroup
  \small
  \setstretch{1.0}
  \setlength{\tabcolsep}{4pt}
  \renewcommand{\arraystretch}{1.18}
  \caption{Selected measurement and timing requirements in Chapter 5.}
  \label{tab:appendix-task-criteria}
  \begin{tabularx}{\textwidth}{@{}>{\raggedright\arraybackslash}p{35mm}
    >{\raggedright\arraybackslash}X@{}}
    \toprule
    \textbf{Task or test} & \textbf{Criterion and interpretation} \\
    \midrule
    Panda capability positioning and waypoints &
    The chapter applies a common 20~mm distance bound in its revised
    assessment. The original generated limits were 50 and 60~mm,
    respectively. \\
    Panda downstream manipulation &
    Terminal error limits are 50~mm for reaching and pushing, 70~mm for
    pick-and-place and dial turning, and 30~mm for drawer opening.
    Object and fixture tasks also require the specified contact and motion
    sequence; pick-and-place requires release and support for at least 0.2~s. \\
    LEAP all-fingertip reach &
    Concatenated Cartesian error across all four fingertips below
    8.94~mm at 2~s. A shorter record cannot establish this endpoint. \\
    LEAP joint-pose matching &
    Maximum joint error below $10^\circ$ at 4~s.
    This whole-hand task differs from the two-joint capability tests. \\
    LEAP object holding &
    More than five control samples within 10~mm of the target and no
    drop samples beyond 300~mm, over the complete 1.5~s interval. \\
    LEAP block manipulation and valve turning &
    At 4~s, block-position error must be below 10~mm and orientation
    error below 0.10~rad; the valve task requires wrapped angle error
    below 0.10~rad. \\
    Skydio capability positioning and waypoints &
    The chapter uses a revised common distance bound of 25~mm.
    These capability tests are separate from full flight tasks. \\
    Skydio point-of-interest orbit &
    The reported retrospective assessment requires a complete circuit.
    Radius and height errors are bounded by 100 and 50~mm, respectively;
    heading towards the centre has a root-mean-square error limit of
    $12^\circ$. \\
    \bottomrule
  \end{tabularx}
  \par\smallskip
  \noindent\raggedright\emph{Notes:} The table summarises the requirements
  relevant to the chapter's interpretation; it is not a replacement for the
  complete executable task definitions. The stored fingertip tolerance is
  0.00894427191~m. Revised criteria and separate retries are distinguished
  in Section~\ref{sec:appendix-reviewed-results}.\par
  \endgroup
\end{table}

\section{Reviewed Results and Separate Executions}
\label{sec:appendix-reviewed-results}

\subsection{Measurement Review in Experiment 2a}

The \emph{Automatic} and \emph{Incl.\ review} columns in
Table~\ref{tab:chapter4-exp2a-generation-results} describe different
assessments of the recorded synthesis attempts.
The automatic column uses the original fixed behavioural test suite for each robot.
Measurement review re-examined selected trajectories where the observation
binding or reference state did not express the intended requirement.
It did not constitute a new driver generation or task execution.

\begin{table}[htbp]
  \centering
  \begingroup
  \small
  \setstretch{1.0}
  \setlength{\tabcolsep}{4pt}
  \renewcommand{\arraystretch}{1.18}
  \caption{Measurement reviews underlying the Experiment 2a results.}
  \label{tab:appendix-exp2a-review}
  \begin{tabularx}{\textwidth}{@{}>{\raggedright\arraybackslash}p{34mm}
    >{\raggedright\arraybackslash}X@{}}
    \toprule
    \textbf{SO-101 runs} & \textbf{Review and treatment} \\
    \midrule
    Sonnet 4.6, runs 01 and 02 &
    The original contact binding named a jaw geometry that did not carry
    the observed contact. Review used the contacting gripper geometry,
    preserving the trajectory, threshold, window and aggregation.
    Both runs contributed additional reviewed passes. \\
    DeepSeek V4 Pro, run 02 &
    The original hold check measured displacement from the action's initial
    pose. Review assessed the post-convergence hold over the original
    1.5--3~s window, retaining the 0.05~rad tolerance.
    The run contributed one additional reviewed pass. \\
    Opus 5, run 03 &
    Correcting the hold reference satisfied that test, but the measured
    contact force was 33.43~N against a requested upper bound of 30~N.
    The original automatic criterion allowed 40~N.
    The run contributed no additional reviewed pass. \\
    \bottomrule
  \end{tabularx}
  \endgroup
\end{table}

The review decisions and links to their underlying measurements are retained
in \path{outcome_review.json} under
\path{expriment/chapter4_synthesis/data/runs/exp2a_8000_01/}.
The original reports remain unchanged.
These measurement reviews are reported separately from the original
automatic assessments.

\subsection{Cross-Robot Reassessment and Retries}
\label{sec:appendix-cross-robot-review}

Chapter 5 uses revised distance tolerances for Panda and Skydio.
For Panda, the positioning and waypoint observations satisfy both their
original limits and the revised 20~mm bound.
For Skydio, the recorded positioning, waypoint and terminal-position errors
also satisfy the revised 25~mm bound. These retrospective comparisons use
existing measurements.

The Skydio task result reported in Chapter 5 is a retrospective 5/5
assessment. The original task reports recorded 3/5 under their timing
and completion checks. Stationary holding exceeded the 60~s limit by
0.01~s. The first orbit continued beyond 60~s, and its terminal net
angular progress fell just short of a full circle despite an earlier
complete circuit. The chapter's retrospective result is therefore distinct
from the original deadline-constrained score; tightening geometric tolerances
alone does not account for the difference.

Retries involved new executions and are reported separately.
The LEAP retry reused the same driver and passed fingertip reaching once
the record covered the required 2~s endpoint. The Skydio orbit retry
completed a circuit but exceeded the heading-error limit.
Neither retry replaces the first attempt or adds an independent driver
generation.

Panda used a separate export of its partially validated driver.
LEAP required missing forwarding wrappers to be completed after automatic
export failed. Skydio underwent separate validation and wrapper preparation
after generation reached its turn limit.
These interventions enabled downstream assessment without changing the
original pipeline outcomes. The records are retained under
\path{expriment/chapter5_cross_robot/data/runs/}.

\section{Model Identities and Cost Assumptions}
\label{sec:appendix-models}

Model names in Chapter 4 denote the endpoint configurations in
Table~\ref{tab:appendix-model-identities}.
The seven-model list describes the benchmark configurations; it does not
imply that all seven yielded comparable synthesis results.
The response limit for the synthesis batch reported in Experiment 2a was
8,000 tokens per model turn. This differs from the model-specific limits
recorded for the earlier seven-model configuration.

\begin{table}[htbp]
  \centering
  \begingroup
  \small
  \setstretch{1.0}
  \setlength{\tabcolsep}{4pt}
  \renewcommand{\arraystretch}{1.18}
  \caption{Endpoint identities and recorded seven-model output limits.}
  \label{tab:appendix-model-identities}
  \begin{tabularx}{\textwidth}{@{}>{\raggedright\arraybackslash}p{31mm}
    >{\raggedright\arraybackslash}X
    >{\raggedleft\arraybackslash}p{18mm}@{}}
    \toprule
    \textbf{Model} & \textbf{Runtime identity} & \textbf{Output cap} \\
    \midrule
    Claude Sonnet 4.6 & \path{eu.anthropic.claude-sonnet-4-6} & 8,192 \\
    Claude Opus 5 & \path{eu.anthropic.claude-opus-5} & 32,768 \\
    Claude Haiku 4.5 & \path{eu.anthropic.claude-haiku-4-5-20251001-v1:0} & 8,192 \\
    Amazon Nova Pro & \path{amazon.nova-pro-v1:0} & 5,000 \\
    DeepSeek V4 Pro & \path{deepseek-v4-pro} & 16,384 \\
    Ministral 3 8B & \path{mistral.ministral-3-8b-instruct} & 8,000 \\
    GPT-5.6 Sol & \path{openai.gpt-5.6-sol} & 32,768 \\
    \bottomrule
  \end{tabularx}
  \par\smallskip
  \noindent\raggedright\emph{Notes:} Output caps are configured tokens per
  response, not total episode budgets or model maximum capacities.
  They reproduce the earlier recorded configuration and do not apply
  uniformly to the later synthesis runs.\par
  \endgroup
\end{table}

Table~\ref{tab:appendix-model-pricing} retains the dated rate assumptions
associated with the benchmark's cost reporting.
For separately recorded uncached input, cached input and output counts,
the token-cost calculation is
\begin{equation}
  C = \frac{N_{\mathrm{in}}p_{\mathrm{in}}
        +N_{\mathrm{cache}}p_{\mathrm{cache}}
        +N_{\mathrm{out}}p_{\mathrm{out}}}{10^6},
\end{equation}
where each $p$ is a dollar rate per million tokens and $C$ is in US dollars.
The input categories must be disjoint to avoid counting cached tokens twice.
Estimates exclude local computation and do not represent provider invoices.

\begin{table}[H]
  \centering
  \begingroup
  \small
  \setstretch{1.0}
  \setlength{\tabcolsep}{4pt}
  \renewcommand{\arraystretch}{1.18}
  \caption{Recorded token-price assumptions, in US dollars per million tokens.}
  \label{tab:appendix-model-pricing}
  \begin{tabularx}{\textwidth}{@{}>{\raggedright\arraybackslash}X
    >{\raggedleft\arraybackslash}p{26mm}
    >{\raggedleft\arraybackslash}p{26mm}
    >{\raggedleft\arraybackslash}p{26mm}@{}}
    \toprule
    \textbf{Model} & \textbf{Uncached input} &
    \textbf{Cached input} & \textbf{Output} \\
    \midrule
    Claude Sonnet 4.6 & 3.30 & 0.33 & 16.50 \\
    Claude Opus 5 & 5.50 & 0.55 & 27.50 \\
    Claude Haiku 4.5 & 1.10 & 0.11 & 5.50 \\
    Amazon Nova Pro & 1.13 & -- & 4.52 \\
    DeepSeek V4 Pro & 0.66 / 1.32 & 0.022 / 0.044 & 1.98 / 3.96 \\
    Ministral 3 8B & 0.23 & -- & 0.23 \\
    GPT-5.6 Sol & 4.00 / 8.00 & 0.40 / 0.80 & 20.00 / 30.00 \\
    \bottomrule
  \end{tabularx}
  \par\smallskip
  \noindent\raggedright\emph{Notes:} A dash denotes an unavailable separate
  cache rate. DeepSeek values are off-peak / peak; GPT-5.6 Sol values are
  standard / long-context. The historical reference date is
  1 September 2026. Rates are retained as recorded assumptions rather
  than updated to current prices.\par
  \endgroup
\end{table}

The Claude rates include the recorded 10\% EU regional premium
\citep{anthropic2026pricing}. The Nova Pro and Ministral assumptions refer
to Amazon Bedrock \citep{aws2026bedrockpricing}.
The exact historical Nova Pro rates could not be independently reconstructed
from the later public pricing page and remain a cost-estimation assumption.
The DeepSeek schedule distinguishes off-peak and peak rates
\citep{deepseek2026pricing,deepseek2026v4pro}.
For GPT-5.6 Sol, the recorded long-context rates apply to the entire request
when input exceeds 272,000 tokens \citep{openai2026gpt56sol}.

The selected synthesis-run costs in
Figure~\ref{fig:chapter4-exp2a-synthesis-traces} use the pricing and token
accounting recorded with that figure.
All reported Claude prompt tokens were charged at the uncached input rate
because the gateway did not provide a cache breakdown.
DeepSeek estimates used off-peak rates and separately recorded cache reads.
These assumptions and the selected run identities are retained in
\path{exp2a_successful_run_turns.json}, under
\path{expriment/chapter4_synthesis/figures/}.
